% Chapter Template

% Main chapter title
%\chapter[toc version]{doc version}
\chapter{Introduction}

% Short version of the title for the header
%\chaptermark{version for header}

% Chapter Label
% For referencing this chapter elsewhere, use \ref{ChapterTemplate}
\label{Introduction}

% Write text in here
% Use \subsection and \subsubsection to organize text
%
%In the modern era of digital communication, livestreaming has become a cornerstone of real-time multimedia delivery, enabling users to share and consume content instantly, from entertainment platforms to mission-critical applications, the demand for seamless and secure livestreaming continues to grow. However, high-concentration areas, such as stadiums, conferences, and urban centers, often face severe bandwidth restrictions due to the high volume of simultaneous users. The rapid growth of livestreaming services has created unprecedented demands on network infrastructure, these areas often experience bandwidth bottlenecks that degrade the quality of service, leading to interruptions, delays, and reduced user satisfaction. Cellular networks and centralized cloud architectures, while widely used, are often too expensive or fail to scale effectively under heavy loads, so Wi-Fi networks offer a cost-effective alternative, but their potential remains underutilized due to limitations in bandwidth allocation and security protocols.
%
%The proliferation of Internet of Things (IoT) devices and edge computing has revolutionized the way data is processed and transmitted, traditional cloud-based architectures, while powerful, often struggle to meet the latency and bandwidth requirements of modern livestreaming applications, Fog computing offers a promising solution by decentralizing computational resources and bringing them closer to the end-users by leveraging fog nodes in conjunction with Wi-Fi networks, this way livestreaming systems can achieve faster data processing, reduced latency, and improved scalability. Despite these advancements, security remains a critical concern, this solutions remain vulnerable to various threats, including data breaches, unauthorized access, and man-in-the-middle attacks, ensuring the confidentiality, integrity, and availability of the data is paramount, especially in applications involving sensitive information. The intersection of fog computing, Wi-Fi optimization, and secure livestreaming presents both opportunities and challenges that require innovative solutions.
%
%To address these challenges, this thesis proposes a comprehensive framework for secure livestreaming with fog offloading over Wi-Fi networks. The approach involves designing and implementing a decentralized security model that leverages the computational capabilities of fog nodes while ensuring end-to-end encryption and authentication, by offloading computational tasks to fog nodes and optimizing Wi-Fi bandwidth allocation, the framework aims to reduce latency and optimize resource utilization without compromising security.
%
%The research methodology combines theoretical analysis, system design, and experimental evaluation to achieve several key objectives: identifying security vulnerabilities in existing livestreaming systems, developing novel algorithms for secure fog offloading over Wi-Fi, and validating the framework through real-world use cases. The ultimate goal is to provide a scalable and cost-effective solution that addresses the deficiencies of current technologies and establishes a foundation for future advancements in secure livestreaming.

In the modern era of digital communication, livestreaming has emerged as a cornerstone of real-time multimedia delivery. It enables users to share and consume content instantaneously, spanning applications from entertainment and gaming to education, telemedicine, and emergency response \cite{livestreaming_use_cases}. As the demand for seamless and secure livestreaming continues to grow, particularly in high-density environments, traditional network infrastructures face significant challenges.

High-concentration areas—such as stadiums, conferences, and urban centers—are particularly prone to bandwidth saturation due to the massive number of simultaneous connections. This congestion often leads to degraded quality of service, including interruptions, delays, and decreased user satisfaction \cite{urban_bandwidth_bottleneck}. Although cellular networks and centralized cloud architectures are widely deployed to support livestreaming, they are frequently either too costly or incapable of scaling effectively under peak loads \cite{cloud_scalability_limitations}. Wi-Fi networks, by contrast, offer a more cost-effective alternative. However, their potential is often underutilized due to inefficient bandwidth allocation mechanisms and legacy security protocols that do not meet the demands of modern applications \cite{wifi_limits_security_bandwidth}.

The proliferation of Internet of Things (IoT) devices and the rise of edge computing paradigms have introduced new opportunities for data processing and transmission. Traditional cloud-based models, while powerful in aggregate computational capacity, often fall short in meeting the latency and bandwidth requirements of real-time applications such as livestreaming \cite{edge_vs_cloud_latency}. Fog computing—a decentralized computing architecture that extends cloud capabilities to the edge of the network—offers a promising solution. By leveraging fog nodes deployed close to end-users, data can be processed with reduced latency, lower bandwidth consumption, and greater scalability \cite{fog_intro, fog_livestreaming_benefits}. When integrated with Wi-Fi infrastructures, fog computing has the potential to transform conventional livestreaming systems into agile, low-latency platforms capable of supporting diverse user demands.

Despite these technological advancements, security remains a critical concern. Livestreaming systems are increasingly targeted by cyber threats, including data breaches, unauthorized access, denial-of-service attacks, and man-in-the-middle interceptions \cite{livestreaming_threats}. In high-stakes environments—such as healthcare, law enforcement, and industrial monitoring—ensuring the confidentiality, integrity, and availability of livestreamed data is non-negotiable \cite{security_critical_apps}. Current security architectures are often insufficient in decentralized environments, highlighting the need for adaptive and robust security models tailored for fog-Wi-Fi ecosystems.

To address these challenges, this thesis proposes a comprehensive framework for secure livestreaming with fog offloading over Wi-Fi networks. The proposed approach involves the design and implementation of a decentralized security model that leverages the computational capabilities of fog nodes while enforcing end-to-end encryption and mutual authentication. Additionally, the framework includes an intelligent bandwidth allocation mechanism optimized for dynamic Wi-Fi environments, enabling efficient resource utilization and low-latency video delivery. By integrating security and performance optimization at the architectural level, this work aims to enhance the reliability, responsiveness, and safety of next-generation livestreaming systems.


